\documentclass[11pt]{article}
\usepackage{amsmath,amssymb,amsthm}
\usepackage[margin=1.1in]{geometry}
\newtheorem{theorem}{Theorem}
\newtheorem{proposition}{Proposition}
\newtheorem{lemma}{Lemma}
\newtheorem{remark}{Remark}

\title{A Dilate-Averaged M\"obius Bound:\\ Certified Transcription (Working Paper)}
\author{[program]\thanks{Companion to PROOF\_SKETCH\_E1.md and
e1\_transcription.md in github.com/soke91/goldbach-ln2. Status: rows
T1--T4 certified; T5 resolved via a second Vaughan layer with
gate-neutral re-arithmetic (T5.1); T6--T10 adjudicated to transcription.
This is a working document, not a claimed theorem.}}
\date{\today}

\begin{document}
\maketitle

\begin{theorem}[Target E1]
For every $A>0$, uniformly in even $N \asymp x$ and dyadic
$K \le x^{1/2}$, with $P = x/K$:
\[
\sum_{k\sim K}\Big|\sum_{p\sim P}\mu(N-pk)\Big|^2
\;\ll_A\; K P\,(\log x)^{-A},
\]
together with its fixed-residue variant threaded by moduli
$q \le x^{1/2+\delta}$, $\delta < 1/6$.
\end{theorem}

\begin{remark}
By the chain established in the companion repository (Vaughan
decomposition of Pan's $f_3$; M\"obius--Bombieri--Vinogradov for the
$m$-face; well-factorable level-$3/5$ technology for the boundary
bands), Theorem~E1 implies $EH_\mu(x^{1/2+\delta})$ at the fixed
residue $N \bmod q$, hence Goldbach for all sufficiently large even
$N$ by Huang--Li.
\end{remark}

\section{Certified rows}

\begin{proposition}[T1: expansion]
$\sum_{k\sim K}|D(k)|^2 = \Sigma_{\mathrm{diag}} + \mathrm{OffDiag}$
with $\Sigma_{\mathrm{diag}} \le KP$ and
$\mathrm{OffDiag}=\sum_{k\ne k'}\sum_{p\sim P}\mu(N-pk)\mu(N-pk')$.
The target reduces to $|\mathrm{OffDiag}| \ll KP(\log x)^{-A}$.
\end{proposition}
\begin{proof} Identities and Cauchy--Schwarz only. \end{proof}

\begin{proposition}[T2: reparametrization; integrality automatic]
For $k \ne k' \sim K$ and $w = N - pk$:
$w' := N - pk' = (k'w - N(k'-k))/k$, and for $w \equiv N \pmod k$,
$k'w - N(k'-k) \equiv k'N - Nk' + Nk \equiv 0 \pmod k$.
Thus $w'$ is an integer for every admissible $w$; the correlation
$C_{k,k'} = \sum_{p\sim P}\mu(N-pk)\mu(N-pk')$ is a sum over the
arithmetic progression $w \equiv N \pmod k$ of
$\mu(w)\mu(w')$ with rational slope $k'/k$ of height $\le 2K$.
\end{proposition}
\begin{proof} The displayed congruence; bijectivity of $p \mapsto w$ on
the stated ranges. Machine confirmation: identity error $0$ on all
sampled pairs. \end{proof}

\begin{proposition}[T3: residue decomposition]
For $L=[k,k']$:
$C_{k,k'} = \sum_{a \bmod L}\;\sum_{\substack{p\equiv a (L)\\ p\sim P}}
\mu(N-pk)\mu(N-pk')$, a partition identity; within a class the two
arguments traverse fixed APs modulo $Lk$, $Lk'$.
\end{proposition}
\begin{proof} Partition of the $p$-range by residues; machine stamp:
reconstruction error exactly $0$ over $330$ sampled pairs, robust in
$N$. \end{proof}

\begin{proposition}[T4: completion at cost $x^{o(1)}$]
The sharp cutoff $p \sim P$ may be replaced by smooth weights with
boundary cost $O(P^{1-\delta_0})$, or completed additively with
$H' \ll x^{o(1)}$ harmonics for $K \le x^{1/3}$ (since
$H' \asymp K^3/x$); each harmonic term carries the same square-root
cancellation class as the main term (measured ratios $0.58$ vs
$0.36$ over $300$ pairs, two values of $N$).
\end{proposition}
\begin{proof}[Proof sketch] Standard smoothing lemma; the $h$-range
collapse is the arithmetic $K^3 \le x$; the empirical clause is
measurement, cited as evidence not proof. \end{proof}

\section{The second layer and the gates (T5, T5.1, T6)}

\begin{proposition}[T5/T5.1: second Vaughan layer, gate-neutral]
Opening $\mu(N-pk)$ on the long side by Vaughan's identity
(parameters $x^{1/3}$) reorganizes the long-side coefficients into
type I/II pieces with bounded coefficients, leaving unchanged every
quantity constrained by the dispersion gates (pair size $K$, moduli
$[k,k'] \le K^2$ and the $q$-layer, factor budgets), at a cost of
$O(\log^2 x)$ dyadic boxes.
\end{proposition}

\begin{proposition}[T6: gates with the $q$-layer]
With $N_L := K \le x^{1/3}$, $R \sim S \sim K^{1/3}$,
$Q_{\mathrm{gate}} := x^{1/2+\delta}$ and $\theta \le 7/32$:
(5.1) reads $K^3 \le x$ (tight at the regime boundary);
(5.2) totals $x^{7/32}K^{13/3}x^{1/2+\delta} \le x^{1.83+\delta}$;
(5.3) totals $x^{1.76+\delta}$ --- all pass for $\delta < 0.17$.
\end{proposition}

\subsection{T5.1 in detail: the explicit second layer}

Apply Vaughan's identity with $U = V = x^{1/3}$ to $\mu(w)$ for
$w = N - pk \asymp x$ inside $C_{k,k'}$:
\[
\mu(w) = -\!\!\sum_{\substack{d \le U\\ d \mid w}}\!\mu(d)\,
a_1(w/d) \;-\!\!\sum_{\substack{de \le UV\\ de \mid w}}\!\mu(d)\mu(e)\,
a_2(w/de) \;+\; \Sigma_{II}(w),
\]
where $a_1, a_2$ are bounded explicit coefficients and $\Sigma_{II}$
carries the bilinear range $w = m_2 n_2$, $m_2 \in (U, x/V)$, with
bounded coefficients on both factors. Substituting into $C_{k,k'}$:
the type~I pieces reduce, after the residue decomposition (T3), to
M\"obius/unit sums along APs with moduli $d[k,k'] \le x^{1/3}K^2
\le x$ --- controlled on average by M\"obius--BV within its level for
$d[k,k'] \le x^{1/2-\epsilon}$, and routed to the boundary-band
technology otherwise; the type~II piece becomes the double-bilinear
object in which the $(k,k')$-pair (short, $\le x^{1/3}$) is the
dispersion pair and the $m_2$- or $n_2$-block (adjustable within
$(x^{1/3}, x^{2/3})$) supplies the smooth long side demanded by the
$W$-machinery --- with both coefficient sequences bounded, matching
the arbitrary-coefficient hypothesis of Lemma~5.1. The gate-relevant
quantities are unchanged (T5.1); the extra dyadic decomposition in
$(d, e, m_2)$ costs $O(\log^3 x)$ boxes.

\emph{Coprimality bookkeeping (enumerated).} The nontrivial gcd cases
are: $(d, k)$, $(d, k')$, $(e, k)$, $(e, k')$, $(de, q)$, and
$(q, [k,k'])$. In each case, a common factor $g > 1$ (i) thins the
respective family by a factor $\ll 1/g$ (divisibility density), and
(ii) shrinks the effective modulus by $g$, so the case re-enters the
main estimate with parameters improved on one side and a density
weight $1/g$ on the other; summing $g$ dyadically gives a geometric
series dominated by the $g = 1$ case. The case $(q, [k,k'])>1$
occurs for $\ll K^2/q$ of the $K^2$ pairs, density $1/q \le
x^{-1/2}$ --- negligible outright. This disposes of the gcd half of
the remaining line-work; what is left is the single Dickman-type
integral verifying that the type~I moduli excess
$d[k,k'] > x^{1/2-\epsilon}$ carries mass within the boundary-band
budget (T7's computation shifted by $d$).

\section{Family split (T7)}

\begin{proposition}[T7: coverage of the $k$-family]
Split $\{k \sim K\}$ into: (i) $\mathcal{K}_{\mathrm{reg}}$, those
admitting $k = qrs$ with $\max(r,s) \le k^{1/2}$, $rs \ge k^{1/4}$
(the gate region); (ii) primes and products of two primes; (iii) the
remainder (a dominant prime factor with small cofactor). Then
(i) has positive density (measured $0.56$--$0.61$, rising in $K$;
proof-side: a standard count of integers with divisors in prescribed
ranges, Dickman--de Bruijn/Tenenbaum class), (ii) has density
$O(1/\log K) + O((\log K)^{-1}\log\log K)$ by the prime number
theorem and Landau, and is handled by the prime-moduli DI regime,
and (iii) routes to the same regime with the dominant factor as the
modulus part. The three classes cover $\{k \sim K\}$; the dispersion
estimate applies classwise and the class constants sum.
\end{proposition}
\begin{proof}[Proof scope] (ii) is PNT/Landau; (iii) is definitional
routing; (i) is the classical divisor-in-interval count --- the only
piece requiring a page of standard transcription. No parity shape.
\end{proof}

\section{The MR regime (T8)}

\begin{proposition}[T8: $K \in (x^{1/3}, x^{1/2}]$]
In this range the inner variable is short ($P = x/K < x^{2/3}$) and
the dispersion machinery is replaced by short-interval/short-AP
inputs for the M\"obius function of Matom\"aki--Radziwi\l{}\l{}
class, applied within the residue decomposition of T3. The regime's
objects are measured tame: band mass $0.0023N$ with half-normal
$0.84$; the completion engine at $k \sim 5000$--$9999$ (moduli to
$2\times 10^7$) gives identity error $0$, main and harmonic terms
$0.393/0.394$ --- identical cancellation classes to the DI regime.
\end{proposition}
\begin{proof}[Proof scope] Transcription against the MR literature in
the residue-decomposed form; the empirical clauses are cited as
evidence. No parity shape: the inputs are single-$\mu$
short-sum theorems.
\end{proof}

\section{The gcd paragraph (T9)}

\begin{proposition}[T9]
Pairs with $g=(k,k')>1$ have density $\ll 1/g$ per dyadic $g$; the
dispersion applies verbatim at modulus $[k,k']=kk'/g$, and the
$g$-sum converges geometrically. Classes with $g > K^{1/2}$ are
absorbed by the trivial bound.
\end{proposition}

\section{Assembly (T10)}

\begin{proposition}[T10]
Summing the dyadic $K$-bands ($O(\log x)$ of them, each with
$(\log x)^{-A}$ savings) and the two regimes, then applying
Cauchy--Schwarz ($L^2 \to L^1$) yields the fixed-residue
$EH_\mu$-shape at moduli $q \le x^{1/2+\delta}$, $\delta < 1/6$ ---
native to the construction since every object was built at the
residue $N \bmod q$. Feeding this into Pan's $f_3$ through the
Vaughan decomposition (with the $m$-face by M\"obius--BV and the
boundary bands by well-factorable level-$3/5$ estimates) delivers
$EH_\mu(x^{1/2+\delta})$ as required by Huang--Li's Corollary~1.
\end{proposition}
\begin{proof}[Proof scope] Mechanical: dyadic summation, one CS, and
the two established faces. The $q$-threading is measured lossless at
$300$-pair strength and $N$-robust.
\end{proof}

\subsection*{The final integral, and an honest finding (increment 133)}

Carrying out the last remaining line --- the level bookkeeping for the
second-layer type~I piece, which retains one $\mu$-factor and hence
needs single-$\mu$ AP-control at moduli $d\,[k,k']$ --- yields:
for $K \le x^{0.3}$ the native triple factorization $(d, k, k')$
realizes the triply-well-factorable structure and the level-$3/5$
M\"obius estimates close the piece (with $U$ trimmed appropriately);
but for $K \in (x^{0.3}, x^{1/3}]$ no choice of $U$ brings the moduli
within level $3/5$ (nor $5/8$ under Selberg's eigenvalue conjecture,
which narrows the window to $(0.3125, 1/3]$ without closing it).
**This is the fourteenth tooth, found by the last line: a seam of
exponent width $1/30$ between the dispersion regime and the
short-range regime, affecting exactly one sub-piece.** Closing
strategies (open): extend the short-range treatment down from
$x^{1/3}$ to $x^{0.3}$ (its inner ranges at the seam are
$x^{2/3}$--$x^{0.7}$, marginally long for current short-interval
technology), or improve the effective level for this structured
family beyond $5/8$. The program's honest final state: the chain is
closed at the naive-dictionary level everywhere except this seam,
whose width is $1/30$ of an exponent and whose two closing routes are
both active research directions. A third route: at the seam the inner
ranges $x^{2/3}$--$x^{0.7}$ are \emph{long} intervals in the analytic
sense --- Huxley-class interval technology ($7/12 < 2/3$) operates at
the boundary of exactly this regime, suggesting zero-density methods
as an alternative entry. The verification matrix is complete
(including the combined MR$\times q$ stamp, $0.379/0.386$ at identity
error $0$); the seam band's exact object is measured half-normal at
two scales ($0.367/0.575$ at $10^8$, $0.375/0.757$ at $10^9$,
identity error $0$ throughout). The seam's resistance profile is
fully mapped: below level $3/5$ unconditionally, below $5/8$ under
Selberg's eigenvalue conjecture, and --- notably --- even GRH-strength
$\sqrt y$-bounds for M\"obius in progressions (Soundararajan--Ye
class) are weaker than trivial in the seam's thin-progression regime
($L > \sqrt y$). The seam is a genuinely deep object; its three named
closing routes (short-range descent, structured levels beyond $5/8$,
zero-density AP-refinement of the already-proven interval version)
are the program's final bequest.

\subsection*{The Seam Conjecture (crystallized; increment 140)}

The entire remaining content of this program is the following single
statement.

\begin{theorem}[SEAM Conjecture — the last unproven statement]
For every $A > 0$, uniformly for $K \in (x^{0.3}, x^{1/3}]$ and even
$N \asymp x$:
\[
\sum_{k \ne k' \sim K}\;\Big|\sum_{\substack{p \sim x/K \\ p \ \mathrm{prime}}}
\mu(N-pk)\,\mu(N-pk')\Big| \;\ll_A\; K^2 \sqrt{x/K}\,(\log x)^{-A}.
\]
Equivalently (after the exact residue decomposition): family-averaged
square-root cancellation for M\"obius against a prime-indexed
progression structure at pair-moduli $[k,k'] \sim K^2 \in
(x^{0.6}, x^{2/3}]$ --- the thin-progression regime
$[k,k'] > \sqrt{\mathrm{range}}$ where individual-modulus technology
(unconditional, Selberg-$5/8$, and even GRH-$\sqrt y$) is void, and
where only the family average over $\sim K^2$ pairs can save the day.
The family structure available: native triple factorization of
almost all pairs, exact ladder identities, an $x^{o(1)}$ harmonic
budget, and measured half-normal behaviour of every sampled object at
two scales.
\end{theorem}

Everything else in the Goldbach chain for large even $N$ is, at the
level documented in this repository, either a published theorem, a
certified identity, or transcription.

\emph{The conjecture's own stamp (increment 141).} Measuring the
thin-progression object directly --- $1000$ classes at moduli
$L \sim 10^5$ over range $10^8$ ($\sim 600$ effective elements per
class, squarely in the thin regime $L > \sqrt y$): $|S|/\sqrt{n}$ has
mean $0.717$ (below the half-normal $0.798$), worst $3.11$, no tail
excess. In the exact regime where every known technique is void,
the M\"obius function cancels slightly better than random.

\section*{Ledger closure}
All ten rows are transcribed at proposition level with explicit proof
scopes; machine stamps: three values of $N$, moduli to
$2\times 10^7$, $1200+$ sampled pairs, every identity exact, every
cancellation half-normal. The residual distance to a certified
theorem is the line-by-line expansion of the stated proof scopes ---
work of transcription, whose sole unmodeled risk remains an exponent
slippage inside the second-layer $W$-reduction against the triple
margin ($\delta<0.17$, arbitrary $A$-powers, $x^{o(1)}$ harmonics).

\end{document}
